% Part: first-order-logic
% Chapter: sequent-calculus
% Section: propositional-rules

\documentclass[../../../include/open-logic-section]{subfiles}

\begin{document}

\olfileid{mvl}{seq}{prl}

\olsection{निवडक तर्कशास्त्रांसाठी विधानीय नियम}

एखाद्या संयोजकासाठीचे $n$-बाजूंच्या क्रमवर्ती कलनातील निगमन नियम
फक्त त्या संयोजकाच्या वैशिष्ट्यपूर्ण सत्यता-फलनावर अवलंबून असतात.
म्हणून भिन्न तर्कशास्त्रांत एखाद्या संयोजकाची व्याख्या एकाच
सत्यता-फलनाने केली असेल, तर त्या संयोजकाचे $n$-बाजूंचे
क्रमवर्ती नियम त्या तर्कशास्त्रांत सारखेच असतात.

\subsection{$\lnot$ साठीचे नियम}

$\lnot$ साठीचे पुढील नियम \L ukasiewicz आणि क्लिनी यांच्या
तर्कशास्त्रांना तसेच त्यांच्या प्रकारांना लागू होतात.

\begin{defish}
  \begin{center}
\AxiomC{$ \Gamma \nSequent \Pi \nSequent \Delta, !A $}
\RightLabel{\iR{\lnot}{\False}}
\UnaryInfC{$ \lnot !A, \Gamma \nSequent \Pi \nSequent \Delta$}
\DisplayProof
\\[2ex]
\AxiomC{$ \Gamma \nSequent !A, \Pi \nSequent \Delta $}
\RightLabel{\iR{\lnot}{\Undef}}
\UnaryInfC{$\Gamma \nSequent \lnot !A, \Pi \nSequent \Delta$}
\DisplayProof
\\[2ex]
\AxiomC{$!A, \Gamma \nSequent \Pi \nSequent \Delta$}
\RightLabel{\iR{\lnot}{\True}}
\UnaryInfC{$\Gamma \nSequent \Pi \nSequent \Delta,  \lnot !A$}
\DisplayProof
  \end{center}
\end{defish}

$\lnot$ साठीचे पुढील नियम G\"odel तर्कशास्त्राला लागू होतात.

\begin{defish}
\AxiomC{$ \Gamma \nSequent !A, \Pi \nSequent \Delta, !A $}
\RightLabel{\iR{\lnot}{\False}[\LogGod]}
\UnaryInfC{$ \lnot !A, \Gamma \nSequent \Pi \nSequent \Delta$}
\DisplayProof
\hfill
\AxiomC{$!A, \Gamma \nSequent \Pi \nSequent \Delta$}
\RightLabel{\iR{\lnot}{\True}[\LogGod]}
\UnaryInfC{$\Gamma \nSequent \Pi \nSequent \Delta,  \lnot !A$}
\DisplayProof
\hfill
\end{defish}

(G\"odel तर्कशास्त्रात $\lnot !A$ ला~$\Undef$ हे मूल्य
कधीही मिळत नाही; म्हणून मधल्या जागेसाठी नियम नाही.)

\subsection{$\land$ साठीचे नियम}

\L ukasiewicz, प्रबळ क्लिनी आणि G\"odel यांच्या
तर्कशास्त्रांतील $\land$ साठीचे नियम पुढीलप्रमाणे आहेत.

\begin{defish}
\begin{center}
  \AxiomC{$!A, !B, \Gamma \nSequent \Pi \nSequent \Delta$}
  \RightLabel{$\iR{\land}{\False}$}
  \UnaryInfC{$!A \land !B, \Gamma \nSequent \Pi \nSequent \Delta$}
  \DisplayProof
\\[2ex]
  \AxiomC{$\Gamma \nSequent !A, \Pi \nSequent !A, \Delta$}
  \AxiomC{$\Gamma \nSequent !B, \Pi \nSequent !B, \Delta$}
  \AxiomC{$\Gamma \nSequent !A, !B, \Pi \nSequent \Delta$}
  \RightLabel{$\iR\land\Undef$}
  \TrinaryInfC{$\Gamma \nSequent !A \land !B, \Pi \nSequent \Delta$}
  \DisplayProof
  \\[2ex]
  \AxiomC{$\Gamma \nSequent \Pi \nSequent \Delta, !A$}
  \AxiomC{$\Gamma \nSequent \Pi \nSequent \Delta, !B$}
  \RightLabel{$\iR\land\True$}
  \BinaryInfC{$\Gamma \nSequent \Pi \nSequent \Delta, !A \land !B$}
  \DisplayProof
\end{center}
\end{defish}

\subsection{$\lor$ साठीचे नियम}

\L ukasiewicz, प्रबळ क्लिनी आणि G\"odel यांच्या
तर्कशास्त्रांतील $\lor$ साठीचे नियम पुढीलप्रमाणे आहेत.

\begin{defish}
\begin{center}
 \AxiomC{$!A, \Gamma \nSequent \Pi \nSequent \Delta$}
  \AxiomC{$!B, \Gamma \nSequent \Pi \nSequent \Delta$}
  \RightLabel{$\iR\lor\False$}
  \BinaryInfC{$!A \lor !B, \Gamma \nSequent \Pi \nSequent \Delta$}
  \DisplayProof
\\[2ex]
  \AxiomC{$!A, \Gamma \nSequent !A, \Pi \nSequent \Delta$}
  \AxiomC{$!B, \Gamma \nSequent !B, \Pi \nSequent \Delta$}
  \AxiomC{$\Gamma \nSequent !A, !B, \Pi \nSequent \Delta$}
  \RightLabel{$\iR\lor\Undef$}
  \TrinaryInfC{$\Gamma \nSequent !A \lor !B, \Pi \nSequent \Delta$}
  \DisplayProof
  \\[2ex]
  \AxiomC{$\Gamma \nSequent \Pi \nSequent \Delta, !A, !B$}
  \RightLabel{$\iR\lor\True$}
  \UnaryInfC{$\Gamma \nSequent \Pi \nSequent \Delta, !A \lor !B$}
  \DisplayProof
\end{center}
\end{defish}

\subsection{$\lif$ साठीचे नियम}

\L ukasiewicz तर्कशास्त्रातील $\lif$ साठीचे नियम पुढीलप्रमाणे आहेत.

\begin{defish}
\begin{center}
  \AxiomC{$\Gamma \nSequent \Pi \nSequent \Delta, !A$}
  \AxiomC{$!B, \Gamma \nSequent \Pi \nSequent \Delta$}
  \RightLabel{$\iR\lif\False[\LogLuk[3]]$}
  \BinaryInfC{$!A \lif !B, \Gamma \nSequent \Pi \nSequent \Delta$}
  \DisplayProof
  \\[2ex]
  \AxiomC{$\Gamma \nSequent !A, !B, \Pi \nSequent \Delta$}
  \AxiomC{$!B, \Gamma \nSequent \Pi \nSequent \Delta, !A$}
  \RightLabel{$\iR\lif\Undef[\LogLuk[3]]$}
  \BinaryInfC{$\Gamma \nSequent !A \lif !B, \Pi \nSequent \Delta$}
  \DisplayProof
\\[2ex]
  \AxiomC{$!A, \Gamma \nSequent !B, \Pi \nSequent \Delta, !B$}
  \AxiomC{$!A, \Gamma \nSequent !A, \Pi \nSequent \Delta, !B$}
  \RightLabel{$\iR{\lif}{\True}[\LogLuk[3]]$}
  \BinaryInfC{$\Gamma \nSequent \Pi \nSequent \Delta, !A \lif !B$}
  \DisplayProof
\end{center}
\end{defish}

प्रबळ क्लिनी तर्कशास्त्रातील $\lif$ साठीचे नियम पुढीलप्रमाणे आहेत.

\begin{defish}
\begin{center}
  \AxiomC{$\Gamma \nSequent \Pi \nSequent \Delta, !A$}
  \AxiomC{$!B, \Gamma \nSequent \Pi \nSequent \Delta$}
  \RightLabel{$\iR\lif\False[\LogKs]$}
  \BinaryInfC{$!A \lif !B, \Gamma \nSequent \Pi \nSequent \Delta$}
  \DisplayProof
  \\[2ex]
  \AxiomC{$!B, \Gamma \nSequent !B, \Pi \nSequent \Delta$}
  \AxiomC{$\Gamma \nSequent !A, !B, \Pi \nSequent \Delta$}
  \AxiomC{$\Gamma \nSequent !A, \Pi \nSequent \Delta, !A$}
  \RightLabel{$\iR\lif\Undef[\LogKs]$}
  \TrinaryInfC{$\Gamma \nSequent !A \lif !B, \Pi \nSequent \Delta$}
  \DisplayProof
\\[2ex]
  \AxiomC{$!A, \Gamma \nSequent \Pi \nSequent \Delta, !B$}
  \RightLabel{$\iR{\lif}{\True}[\LogKs]$}
  \UnaryInfC{$\Gamma \nSequent \Pi \nSequent \Delta, !A \lif !B$}
  \DisplayProof
\end{center}
\end{defish}

G\"odel तर्कशास्त्रातील $\lif$ साठीचे नियम पुढीलप्रमाणे आहेत.

\begin{defish}
\begin{center}
  \AxiomC{$\Gamma \nSequent !A, \Pi \nSequent \Delta, !A$}
  \AxiomC{$!B, \Gamma \nSequent \Pi \nSequent \Delta$}
  \RightLabel{$\iR\lif\False[\LogGod[3]]$}
  \BinaryInfC{$!A \lif !B, \Gamma \nSequent \Pi \nSequent \Delta$}
  \DisplayProof
  \\[2ex]
  \AxiomC{$\Gamma \nSequent !B, \Pi \nSequent \Delta$}
  \AxiomC{$\Gamma \nSequent \Pi \nSequent \Delta, !A$}
  \RightLabel{$\iR\lif\Undef[\LogGod[3]]$}
  \BinaryInfC{$\Gamma \nSequent !A \lif !B, \Pi \nSequent \Delta$}
  \DisplayProof
\\[2ex]
  \AxiomC{$!A, \Gamma \nSequent !B, \Pi \nSequent \Delta, !B$}
  \AxiomC{$!A, \Gamma \nSequent !A, \Pi \nSequent \Delta, !B$}
  \RightLabel{$\iR{\lif}{\True}[\LogGod[3]]$}
  \BinaryInfC{$\Gamma \nSequent \Pi \nSequent \Delta, !A \lif !B$}
  \DisplayProof
\end{center}
\end{defish}


\begin{sidewaysfigure}
  \begin{center}
  \AxiomC{$A \nSequent A \nSequent A$}
  \RightLabel{\iR \Weakening \True}
  \UnaryInfC{$A \nSequent A \nSequent B, A$}
  \RightLabel{\iR \Weakening \Undef}
  \UnaryInfC{$A \nSequent B, A \nSequent B, A$}
  \RightLabel{\iR \Weakening \Undef}
  \UnaryInfC{$A \nSequent A, B, A \nSequent B, A$}
  \AxiomC{$A \nSequent A \nSequent A$}
  \RightLabel{\iR \Weakening \True}
  \UnaryInfC{$A \nSequent A \nSequent A, A$}
  \RightLabel{\iR \Weakening \True}
  \UnaryInfC{$A \nSequent A \nSequent B, A, A$}
  \RightLabel{\iR \Weakening \False}
  \UnaryInfC{$B, A \nSequent A \nSequent B, A, A$}
  \RightLabel{$\iR\lif\Undef$}
  \BinaryInfC{$A \nSequent A \lif B, A \nSequent B, A$}
  \AxiomC{$B \nSequent B \nSequent B$}
  \RightLabel{\iR \Weakening \Undef}
  \UnaryInfC{$B \nSequent A, B \nSequent B$}
  \RightLabel{\iR \Exchange \Undef}
  \UnaryInfC{$B \nSequent B, A \nSequent B$}
  \RightLabel{\iR \Weakening \Undef}
  \UnaryInfC{$B \nSequent A, B, A \nSequent B$}
  \RightLabel{\iR \Weakening \False}
  \UnaryInfC{$A, B \nSequent A, B, A \nSequent B$}
  \RightLabel{\iR \Exchange \False}
  \UnaryInfC{$B, A \nSequent A, B, A \nSequent B$}
  \AxiomC{$A \nSequent A \nSequent A$}
  \RightLabel{\iR \Weakening \True}
  \UnaryInfC{$A \nSequent A \nSequent B, A$}
  \RightLabel{\iR \Weakening \False}
  \UnaryInfC{$B, A \nSequent A \nSequent B, A$}
  \RightLabel{\iR \Weakening \False}
  \UnaryInfC{$B, B, A \nSequent A \nSequent B, A$}
  \RightLabel{$\iR\lif\Undef$}
  \BinaryInfC{$B, A \nSequent A \lif B, A \nSequent B$}
  \RightLabel{$\iR\lif\False$}
  \BinaryInfC{$A \lif B, A \nSequent A \lif B, A \nSequent B$}
  \DisplayProof
  \end{center}
  \caption{$\LogLuk[3]$ मधील !!{derivation} चे उदाहरण}
\end{sidewaysfigure}

\end{document}
